% Generated draft: requires mathematical and source review.
\documentclass[11pt]{article}
\usepackage[margin=25mm]{geometry}
\usepackage{fontspec}
\setmainfont{Latin Modern Roman}
\usepackage{amsmath,amssymb,mathtools,unicode-math}
\setmathfont{Latin Modern Math}
\usepackage{enumitem,needspace,xurl,hyperref}
\setlength{\parindent}{0pt}
\setlength{\parskip}{5pt}
\setlength{\emergencystretch}{3em}
\newcommand{\N}{\mathbb{N}}
\newcommand{\Z}{\mathbb{Z}}
\newcommand{\Q}{\mathbb{Q}}
\newcommand{\R}{\mathbb{R}}
\newcommand{\C}{\mathbb{C}}
\newcommand{\F}{\mathbb{F}}
\newcommand{\A}{\mathbb{A}}
\newcommand{\PP}{\mathbb{P}}
\newcommand{\E}{\mathbb{E}}
\newcommand{\eps}{\varepsilon}
\newcommand{\dd}{\,\mathrm{d}}
\begin{document}
\section*{UnsolvedMath ID: 9300003}
\textbf{Source problem code:} catalogue entry
\subsection{Definitions and mathematical statement}
Fermat's Last Theorem is the assertion
\[
 \forall n\in\mathbb Z_{\ge3}\ \forall a,b,c\in\mathbb Z_{>0}:\qquad
 a^n+b^n\ne c^n.
\]
The source's question is not whether this assertion is true: it asks whether there is a proof using only mathematical techniques available in the seventeenth century.

That restriction is historical rather than a specified formal proof system. To turn it into a proof-existence statement one would need an exact collection of allowed concepts, axioms, inference methods and historical availability criteria. If such a system \(\mathcal T_{17}\) were fixed, one could ask whether it admits a finite proof of the displayed theorem, but the source supplies no such system. ``Elementary'' is also not automatically equivalent to ``available in the seventeenth century''. No independence or proof-complexity formulation is substituted.

The supplied summary reports the underlying theorem proved by Wiles and Taylor--Wiles in 1994--1995, with the gap in the initial announcement repaired, and separately mentions the later full modularity work of Breuil--Conrad--Diamond--Taylor (2001). That later development is not treated as the date when Fermat's theorem first became proved. The summary says no restricted proof or impossibility result is known and calls the request historical/folklore; its judgment of plausibility is not a mathematical obstruction. The registry's open label applies only to this undefined restriction, not to Fermat's Last Theorem itself.
\subsection{Short English statement}
Can the already-proved Fermat's Last Theorem be proved using only genuinely seventeenth-century methods? The allowed methods are not defined as a formal system, so this is a historical proof-search question rather than a precisely stated independence problem.
\subsection{Sources}
\begin{itemize}
\item Local registry: \url{data/unsolvedmath/problems.json}, permanent UnsolvedMath ID \(9300003\), source code catalogue entry; supplied statement, background, and research summary, updated 2026-08-06.
\item \emph{Unsolved problems in number theory, logic, and cryptography} (2021), FLT.htm: \url{https://unsolvedproblems.org/index_files/FLT.htm}. HTML-text extraction accessed 2026-07-29. The supplied provenance reports inspection of 24 home-page problem links and exclusion of duplicates and source-resolved/nonmathematical items.
\item Wiles and Taylor--Wiles (1994--1995), and Breuil--Conrad--Diamond--Taylor (2001), as named by the summary; no full original-paper citations are supplied.
\end{itemize}
\end{document}
